\documentclass[a4paper, 11pt]{article} % das Papierformat zuerst
\usepackage[utf8]{inputenc}
\usepackage{geometry}
\geometry{a4paper, top=25mm, left=35mm, right=25mm, bottom=25mm}
\usepackage{graphicx}
\usepackage{color}
\usepackage{epstopdf}
\usepackage[T1]{fontenc}
\usepackage{setspace}
\usepackage{tabularx}
\usepackage{blindtext}
\usepackage[ngerman]{babel} %deutsche Silbentrennung
\usepackage{titlesec}
\usepackage{enumitem} 
\usepackage{ifthen}
\usepackage[colorlinks=true,linkcolor=fsi]{hyperref}

\definecolor{fsi}{RGB}{0,81,150}

\newcounter{summe}

\newcommand{\abstimmung}[4]{
	\setcounter{summe}{#3}
	\addtocounter{summe}{#4}
	\begin{flushleft}
		#1\\
	Es wird über den Antrag abgestimmt.
	\end{flushleft}
	\ifthenelse{\equal{#3}{0}\AND\equal{#4}{0}}{
	\begin{center}
		\textbf{Der Antrag wird einstimmig angenommen.}
	\end{center}
	}{
	\begin{center}
		#2 \ifthenelse{\equal{#2}{1}}{Stimme}{Stimmen} dafür, #3 \ifthenelse{\equal{#3}{1}}{Stimme}{Stimmen} dagegen, #4 \ifthenelse{\equal{#42}{1}}{Enthaltung}{Enthaltungen}\\
		\ifthenelse{#2>\value{summe}}{
		\textbf{Der Antrag ist somit angenommen.}
		}{
		\textbf{Der Antrag ist somit abgelehnt.}
		}
	\end{center}
	}  
}


% VARIABLEN
\newcommand{\protokoller}{Max Mustermann}
\newcommand{\dateOfMeeting}{22. Juni 2015}
\newcommand{\TeXer}{Martina Mustermann}
\newcommand{\fsiPresident}{Georg Schäfer}
\newcommand{\protocolID}{8}

\begin{document}
%deckblatt start

\doublespacing
\thispagestyle{empty}

\begin{center}
\includegraphics[width=10.0cm]{../logo_fsi.eps}

\vspace*{\fill}
{\LARGE \textbf{Protokoll der Sitzung des Fachschaftsrates \\vom \dateOfMeeting}}\\
Fachschaftsrat Informatik\\
Trier University of Applied Sciences\\
\vspace{2.5cm}
\textit{
	Protokollführer: \textbf{\protokoller} \\
	\LaTeX - Umsetzung von \TeXer\\
	am \today\\
}
Protokoll ID: \textbf{\#\protocolID}
\vfill
\end{center}

\hspace*{-35cm}
\textcolor{fsi}{\rule{64.9cm}{15pt}}
\pagebreak
 
\setcounter{tocdepth}{2}
\tableofcontents 
\pagebreak

\section{Eröffnung}
Als Protokollführer wird \protokoller~bestimmt.\\
Der Sprecher des Fachschaftsrates \fsiPresident~eröffnet die Sitzung um 19:23 Uhr.
\\\\
\textbf{Es wird festgestellt, dass der Fachschaftsrat beschlussfähig ist.}\\
\textbf{Es fehlen: Santo Pfingsten}
\\\\
Das Protokoll der letzten Sitzung wurde verlesen und genehmigt.
%Die Genehmigung des Protokolls der letzten Sitzung wird auf die nächste Sitzung vertagt.

%\abstimmung{Es wird der Antrag gestellt, Dinge zu tun}{5}{6}{2} -- {dafür}{dagegen}{enthaltungen}

\section{Events}
gdw: 3 leute fachseminar, raeume reserviert, termin festgelegt
100 tage: fabschlussgrillen: tobi und santo treffen sich mit janine.
fabio ist groessenwahnsinnig, will aula fuer notld
notl soll noch vor reading weak veranstaltet werden
spieleabend: am 09.07, noch spiele da, derenbesitzer unbekannt

\section{NIV}
geld fuer server endlich ueberwiesen. 
banner kommt morgen. moglicherweise rabatt beantragen, da verspaetete lieferung.
antraege fuer kabel genehmigt und bestellt.
34 kaesten mate bestellt.
medienanlage: wird bei aufbau geklaert.
uni konnte plakate nicht aufhaengen
bis jetzt 51 reservierungen
auf-/abbau,verkauf: bisher wenige zusagen in doodle

\section{Finanzen}
wir heben kassenbuch. mit stempel. in ablage liste fuer anschaffungen, bald auch auf webseite formular.
noetig: batterien
kasten mate verschwunden
schrank und sofa noch nix neues, dominik schickt moegloches angebot fuer schrank an rat.
alle antraege, die bei asta lagen sind genehmigt
marco schlaegt vor, beamer zu kaufen. wir koennten oechsle / fachbereich fragen.

\section{Raum}
nicht-rats-mitglieder haben noch keinen raumzugang.
putzplan klappt nicht, schwer zu kontrollieren.
wir haben neue tische.
pflanze lebt noch. Warum?!

\section{Sonstiges}
\abstimmung{Es wird der Antrag gestellt, ab jetzt die Beschluesse durchzunummerieren}{12}{0}{0}

Werbung fuer Fachschftsrat Info der Uni wird weiterhin verteilt. Georg setzt sich stark dafuer ein. Stefan ebenso.

\abstimmung{Es wird der Antrag gestellt, keine Werbung mehr fuer den Fachschaftsrat Informatik an der Hochschule zu machen}{2}{6}{4}

\fsiPresident~schliesst sitzung um 19:57

\pagebreak
\section{Bestätigung des Protokolls}
Der Sprecher des Fachschaftsrates Informatik sowie der Protokollführer dieses Protokolls bestätigen mit Ihrer Unterschrift unter diesem Protokoll, dass selbiges inhaltlich korrekt ist und alle darin aufgeführten Beschlüsse so wie beschrieben vom Rat getragen werden.
\\

\vspace{3.5cm}
\hrulefill \hfill \hrulefill

\fsiPresident \hfill \protokoller

{\footnotesize (Sprecher)\hfill (Protokollführer)}
\end{document}
